% talk.tex -- 40-minute session deck for
% "Competition Creates Competition: Stock Prices and the Discovery of Takeover Bidders" (Austin Li).
% Contract: talk/structure-plan.md (sections 1-11) and talk/notes/audit-adoptions.md.
% Every number on a slide carries a "% source:" comment naming its registry key in
% numerics/quantity_registry.csv, or the paper table file (tables/table*.tex) it comes from.
% Frame labels: F0-F16 (F10a, F10b) for the main deck, A1-A29 for the appendix (plan ids).
% Navigation: main frames carry \hypertarget{main:...}; backups carry \hypertarget{app:...}
% and a \BackButton to their single origin. No overlays: the one build (F10a -> F10b) is a
% duplicated frame.
% Look: moloch (maintained metropolis successor) in the UCL brand palette, Fira Sans text at 10 pt
% (reads at about the size of the previous 11 pt Latin Modern Sans), serif math (F29), XeLaTeX.
\documentclass[10pt,aspectratio=169]{beamer}
\usetheme[progressbar=frametitle,
  frametitle margin left=10.95pt,frametitle margin right=10.95pt]{moloch}
\setbeamertemplate{page number in head/foot}[totalframenumber]
% Keep the Madrid text width (433 pt) the frames were laid out for; moloch's default 1 cm margins
% leave 398 pt. Frame titles start at the same left edge as the body.
\setbeamersize{text margin left=10.95pt,text margin right=10.95pt}
\usepackage{fontspec}
% Fira Sans by file name, so the deck compiles wherever TeX Live's fira package is installed
\setsansfont{FiraSans}[Extension=.otf, UprightFont=*-Regular, BoldFont=*-SemiBold,
  ItalicFont=*-Italic, BoldItalicFont=*-SemiBoldItalic]
% UCL brand palette: dark purple leads; bright and mid purple are accents
\definecolor{uclDark}{HTML}{361A54}
\definecolor{uclBright}{HTML}{993BFF}
\definecolor{uclMid}{HTML}{BA82FF}
\setbeamercolor{normal text}{fg=black}
\setbeamercolor{structure}{fg=uclDark}
\setbeamercolor{frametitle}{bg=uclDark,fg=white}
\setbeamercolor{progress bar}{fg=uclBright,bg=uclDark!60}
\setbeamerfont{framesubtitle}{series=\mdseries,size=\small}
\usepackage{econ-slides-compat}
% Moloch sizes its progress bar by \insertframenumber, which the adapter turns into "A1, A2, ..."
% in the appendix. Drive the bar from the frame counter instead and leave it out of the appendix.
\makeatletter
\newif\iftalk@appendix
\let\talk@AppendixStart\AppendixStart
\renewcommand{\AppendixStart}{\global\talk@appendixtrue\talk@AppendixStart}
\setbeamertemplate{progress bar in head/foot}{%
  \nointerlineskip
  \iftalk@appendix\else
  \pgfmathsetlength{\moloch@progressinheadfoot}{%
    \paperwidth*min(1,\number\value{framenumber}/max(1,\inserttotalframenumber))}%
  \begin{beamercolorbox}[wd=\paperwidth]{progress bar in head/foot}
    \begin{tikzpicture}
      \fill[bg] (0,0) rectangle (\paperwidth,\moloch@progressinheadfoot@linewidth);
      \fill[fg] (0,0) rectangle (\moloch@progressinheadfoot,\moloch@progressinheadfoot@linewidth);
    \end{tikzpicture}
  \end{beamercolorbox}\fi}
\makeatother
% The adapter captures the structure color with an ungrouped \usebeamercolor{structure} at
% \begin{document}; that leaks structure purple into every itemize body. Re-select the normal
% text color afterwards (hooks run in order), so list text stays black.
\AtBeginDocument{\usebeamercolor{normal text}}
\usefonttheme[onlymath]{serif}
\usepackage{amsmath,amssymb,mathtools}
\usepackage{booktabs}
\usepackage{array}
\usepackage{tikz}
\usetikzlibrary{arrows.meta,positioning}
\graphicspath{{figures/}}

% COLOR LEDGER (structure-plan section 5; one color per concept, held for the whole talk):
%   UCL purple = chrome only (title bar, progress bar, cover, \KeyIdea, ResultBox, buttons)
%   cInfo = target-payoff spread Delta_T, the information force, "trading-cost window", the chain
%   cCost = preparation costs C, c_L, c_H, rho; "low-cost floor", "high-cost window"
%   deterrence objects (g_H, g_L, B_r) = neutral black; status labels = gray \scriptsize
\colorlet{cInfo}{cAccentB!80!black}  % #AA4B00, contrast 5.65:1 on white: information force (new)
\colorlet{cCost}{cAccentC!70!black}  % #006F50, contrast 6.20:1 on white: preparation-cost objects (F30)

% Local helpers (no new look; gray bookkeeping and navigation only)
% Gray bookkeeping text is a darker gray (#616161, about 6:1 on white) so it survives projection.
\colorlet{cGrayText}{black!62}
\newcommand{\Status}[1]{{\color{cGrayText}\scriptsize #1}}
\newcommand{\GrayLine}[1]{\par{\color{cGrayText}\footnotesize #1\par}}
% Gray inline citations: defined here as well, so the deck does not depend on which copy of the
% adapter is on TEXINPUTS; footnotesize and the darker gray for projection.
\providecommand{\graycite}{}
\renewcommand{\graycite}[1]{{\color{cGrayText}\footnotesize #1}}
\newcommand{\NavSep}{\hspace{0.35em}}
% Navigation buttons in structure purple with white text (beamer's pale default is about 2.5:1)
\setbeamercolor{button}{use=structure,bg=structure.fg,fg=white}
\setbeamercolor{button border}{use=structure,fg=structure.fg}
% Status text readable from the back of the room, for status cells inside main-frame tables
\newcommand{\StatusF}[1]{{\color{cGrayText}\footnotesize #1}}
\newcommand{\Info}[1]{\textcolor{cInfo}{#1}}
\newcommand{\Cost}[1]{\textcolor{cCost}{#1}}
\newcommand{\PropItem}[1]{\par\noindent\hangindent=1.9em\hangafter=1\makebox[1.9em][l]{#1}}% labelled item inside the ResultBox
% Navigation row: the buttons sit in their own band just above the footline (0.5 cm up instead of
% the adapter's 0.75 cm), and frames with buttons reserve that band below their last body line, so
% no text line shares a baseline band with the buttons.
\renewcommand{\PlaceNavShift}[2]{%
  \begin{tikzpicture}[remember picture,overlay]
    \node[anchor=south east, inner sep=0pt, outer sep=0pt]
      at ([xshift=#1,yshift=0.5cm]current page.south east) {#2};
  \end{tikzpicture}}
% Takeaway line with a hanging indent: wrapped lines align with the text after the bullet.
\newcommand{\TakeawayBody}[1]{\begin{minipage}{\textwidth}\raggedright\normalsize
    \setbox0=\hbox{{\usebeamercolor[fg]{structure}\raise0.7pt\hbox{\scriptsize$\bullet$}}\hspace{0.55em}}%
    \hangindent=\wd0 \hangafter=1 \noindent\box0 #1\par
  \end{minipage}}
\renewcommand{\Takeaway}[2][\ExhibitReadingGap]{\par\vspace{#1}\noindent\TakeawayBody{#2}\par}
\renewcommand{\TakeawayWithNav}[2][\ExhibitReadingGap]{\par\vspace{#1}\noindent\TakeawayBody{#2}\par\vspace{0.45cm}}
% Frame 4 timeline box: every box the same fixed height, text top-aligned and ragged right
\newcommand{\ModelBox}[2]{\parbox[t][2.3cm][t]{2.45cm}{\footnotesize\raggedright\hyphenpenalty=10000\relax
  \hspace*{7pt}\textbf{#1}\par\vspace{1pt}#2}}
\newcommand{\NavGap}{\par\vspace{0.45cm}}% bottom reserve on a frame with buttons and no takeaway
\let\AdapterBackButton\BackButton
\renewcommand{\BackButton}[1]{\NavGap\AdapterBackButton{#1}}% a Back button reserves the band too

\title[Competition Creates Competition]{Competition Creates Competition:\texorpdfstring{\\[2pt]}{ }%
Stock Prices and the Discovery of Takeover Bidders}
\author{Austin Li}
\institute{Department of Economics\\University College London}
\date{}

\begin{document}

% ======================================================================
% Frame 0. Title (plain, uncounted): full-bleed UCL dark purple, white text
% ======================================================================
{\setbeamercolor{background canvas}{bg=uclDark}
\setbeamercolor{normal text}{fg=white}\usebeamercolor[fg]{normal text}
\setbeamercolor{title}{fg=white}\setbeamercolor{author}{fg=white}
\setbeamercolor{institute}{fg=uclMid}\setbeamercolor{title separator}{fg=uclBright}
\begin{frame}[plain,noframenumbering,label=F0]
\titlepage
\end{frame}}

% ======================================================================
% Frame 1. Key question
% ======================================================================
\begin{frame}[label=F1]{Does a stronger incumbent keep the challenger out?}
\framesubtitle{The decision interval this paper is about}
\hypertarget{main:interval}{}\hypertarget{main:evidence}{}
\begin{center}
\begin{tikzpicture}[
  box/.style={draw=black!55, rounded corners=3pt, line width=0.6pt, align=center,
    text width=3.9cm, minimum height=0.85cm, font={\footnotesize\hyphenpenalty=10000\relax}, inner sep=3pt},
  arr/.style={-{Stealth[length=5pt]}, line width=0.7pt, black!60}]
  \node[box] (a) {Approach or review disclosed};
  \node[box, right=0.8cm of a] (b) {Target stock keeps trading};
  \node[box, right=0.8cm of b, draw=cCost, line width=1.1pt] (c)
    {Challenger decides whether to pay to prepare};
  \draw[arr] (a) -- (b);
  \draw[arr] (b) -- (c);
\end{tikzpicture}
\end{center}
\vspace{-6pt}
{\small
\begin{itemize}\setlength{\itemsep}{1pt}
\item The \textbf{incumbent} bidder is already prepared; the market knows how strong it is, not what it will bid
\item A \textbf{challenger} must first pay to prepare an executable bid (diligence, financing, approvals); paying is \textbf{entry}
\item \textbf{Received answer:} a stronger rival lowers what preparation earns, so it deters entry \graycite{Fishman 1988; Hirshleifer and Png 1989}
\item \textbf{But} the challenger decides while the stock trades, and it can watch the price
% source: Imprivata chronology, paper/main_filled.md line 39 (no registry number)
\item[] {\color{cGrayText}\footnotesize Imprivata's 2016 proxy separates an unsolicited approach, outreach to potential buyers, and indications of interest conditional on further diligence: a costly preparation stage, not a price drawing anyone in\par}
\end{itemize}}
\TakeawayWithNav[0.5em]{\KeyIdea{The twist:} what that price reveals depends on how strong the incumbent is}
\PlaceNav{\hyperlink{app:interval}{\beamergotobutton{Which deals fit}}\NavSep\hyperlink{app:evidence}{\beamergotobutton{Evidence: a design}}}
\end{frame}

% ======================================================================
% Frame 2. This paper (punchline)
% ======================================================================
\begin{frame}[label=F2]{This paper}
An informed investor trades the target's stock, which pays off from the sale price, before a challenger decides whether to pay to prepare a takeover bid; prices are rational, and trading and entry are solved in equilibrium.
\vspace{1pt}
\begin{itemize}\setlength{\itemsep}{2pt}
\item \KeyIdea{Competition creates competition}: a stronger incumbent can turn an uninformative stock price into an informative one and raise entry\par
  \vspace{1pt}{\small Benchmark: entry \textbf{0.250 % source: base_entry_weak
  $\boldsymbol{\to}$ 0.523} % source: base_entry_strong
  while the challenger's profit before any news falls \textbf{4.80 % source: base_profit_prior_weak
  $\boldsymbol{\to}$ 4.29}\par}% source: base_profit_prior_strong
\item \textbf{Two possible outcomes of one auction}: a strong incumbent can beat a low-value challenger but not a high-value one, so the challenger keeps less while the sale price depends more on who shows up; informed trading puts that into the stock price, and good news justifies expensive preparation
\item \textbf{Information is the channel:} hold fixed what the price reveals, and a stronger incumbent cannot raise entry, as the textbook says \Status{(sign analytical, Prop.~A.3)}
\end{itemize}
\Takeaway[0.45em]{\textbf{Implication:} when the target trades, a stronger rival is not a pure deterrent, because who competes depends on what the price reveals before anyone prepares.}
\end{frame}

% ======================================================================
% Frame 3. Positioning
% ======================================================================
\begin{frame}[label=F3]{What is new relative to learning from prices}
\framesubtitle{Four closest antecedents}
\hypertarget{main:lit}{}\hypertarget{main:refs}{}
\begin{itemize}\setlength{\itemsep}{8pt}
\item \textbf{Prices guide real decisions} \graycite{Dow, Goldstein and Guembel 2017; Edmans, Goldstein and Jiang 2015}: here the sale rule splits one surplus into a traded claim and the challenger's claim, and competition moves them in opposite directions
\item \textbf{Stronger incumbents deter} \graycite{Fishman 1988}: kept intact; I add a market that trades before preparation
\item \textbf{Auction entry is endogenous} \graycite{Levin and Smith 1994}: here entry responds to what the price reveals
\end{itemize}
\TakeawayWithNav[\ExhibitReadingGapRoomy]{Increment: one sale rule moves the two claims oppositely, producing the entry reversal.}
\PlaceNav{\hyperlink{app:lit}{\beamergotobutton{Related work}}\NavSep\hyperlink{app:refs}{\beamergotobutton{References}}}
\end{frame}

% ======================================================================
% Frame 4. Model (D1 timeline)
% ======================================================================
\begin{frame}[label=F4]{Model: who moves, who knows what, and when}
\framesubtitle{The challenger sees the price and its own cost, never the order flow, $\theta$ or $R$}
\hypertarget{main:omit}{}\hypertarget{main:info}{}\hypertarget{main:eq}{}
\begin{center}
\begin{tikzpicture}[
  box/.style={draw=black!55, rounded corners=3pt, line width=0.6pt, align=left,
    text width=2.45cm, inner xsep=4pt, inner ysep=4pt},
  num/.style={circle, fill=black!65, text=white, font=\scriptsize\bfseries,
    inner sep=0pt, minimum size=11pt},
  arr/.style={-{Stealth[length=5pt]}, line width=0.7pt, black!60}]
  \node[box] (s) {\ModelBox{Seller}{commits to a cash second-price auction, reserve $p$}};
  \node[box, right=0.33cm of s] (i) {\ModelBox{Investor}{knows $\theta$; trades the target's stock against noise traders}};
  \node[box, right=0.33cm of i] (m) {\ModelBox{Market makers}{see only total order flow; price the share at expected sale proceeds}};
  \node[box, right=0.33cm of m, draw=cCost, line width=1.1pt] (c)
    {\ModelBox{Challenger}{sees the price and its cost $C$; entry $=$ pay $C$, learn $\theta$, can bid}};
  \node[box, right=0.33cm of c] (b) {\ModelBox{Bidders}{incumbent value $R\sim U[0,r]$; challenger worth $h$ or $\ell$; bid truthfully}};
  \foreach \x/\n in {s/1, i/2, m/3, c/4, b/5} {\node[num] at ([xshift=2pt,yshift=-1pt]\x.north west) {\n};}
  \draw[arr] (s) -- (i); \draw[arr] (i) -- (m); \draw[arr] (m) -- (c); \draw[arr] (c) -- (b);
\end{tikzpicture}
\end{center}
\vspace{-2pt}
{\small
\begin{itemize}\setlength{\itemsep}{4pt}
\item $0<p<\ell<r<h$ ($p$ = reserve, not the stock price): the incumbent can beat a low-value challenger, never a high-value one. Strength $r$ is public; only the incumbent knows its value $R$. Quality $\theta\in\{H,L\}$ (worth $h$ or $\ell$), equally likely, learned only by preparing.
\item \Cost{\textbf{Cost:} preparation cost $C\in\{c_L,c_H\}$, $\Pr(C=c_L)=\rho$, independent of $\theta$ (cheap / expensive preparation)}
\end{itemize}}
\NavGap
\PlaceNav{\hyperlink{app:omit}{\beamergotobutton{What is left out}}\NavSep\hyperlink{app:info}{\beamergotobutton{Why would the investor know?}}\NavSep\hyperlink{app:eq}{\beamergotobutton{Equilibrium}}}
\end{frame}

% ======================================================================
% Frame 5. One auction, two claims
% ======================================================================
\begin{frame}[label=F5]{One auction, two claims, opposite responses}
\framesubtitle{Cash second-price auction with reserve $p$: who wins and who pays at incumbent value $R$}
\hypertarget{main:payoffs}{}\hypertarget{main:payrule}{}
\vspace{4pt}
\begin{center}\small
\begin{tabular}{@{}llll@{}}
\toprule
Incumbent value & High-value challenger & Low-value challenger & Gap in target proceeds\\
\midrule
$R\le\ell$ & wins, pays $\max\{p,R\}$ & wins, pays $\max\{p,R\}$ & $0$\\
$R>\ell$ & wins, pays $R$ & loses; incumbent pays $\ell$ & $R-\ell$\\
\bottomrule
\end{tabular}
\end{center}
\vspace{-2pt}
{\small As incumbent strength $r$ rises:\par}
\vspace{-8pt}
\begin{align*}
\Info{\Delta_T=t_H-t_L=\mathbb{E}[(R-\ell)_+]}\ &\Info{\uparrow} && \text{\Info{target-payoff spread = expected gap (last column)}}\\
g_H=\mathbb{E}[(h-\max\{p,R\})_+]\ &\downarrow && \text{challenger's gross profit ($g_L$ likewise)}
\end{align*}
\vspace{-6pt}
\GrayLine{$t_\theta$ = expected target proceeds when a $\theta$-challenger enters. Proposition 1 (analytical): both signs hold weakly for any first-order strengthening, $R\sim F$ on $[0,\bar r]$, $\ell<\bar r<h$ (strictly in the uniform benchmark, $\bar r=r$).}
\TakeawayWithNav{One shift, opposite signs: challenger keeps less; target proceeds depend more on who it is.}
\PlaceNav{\hyperlink{app:payoffs}{\beamergotobutton{Closed forms}}\NavSep\hyperlink{app:payrule}{\beamergotobutton{Other payment rules}}}
\end{frame}

% ======================================================================
% Frame 6. X1: spread up, profit down
% ======================================================================
\begin{frame}[label=F6]{A stronger incumbent: spread up, profit down}
\framesubtitle{Declared benchmark; weak $r_0=1.2$, % source: base_r_weak
strong $r_1=3$; % source: base_r_strong
auction-stage payoffs, before trading and entry}
\centering
\includegraphics[width=0.92\linewidth,height=0.6\textheight,keepaspectratio]{two_returns_talk.pdf}
\par\vspace{-2pt}{\raggedleft\Status{analytical (Proposition 1)}\par}
\Takeaway{From $r_0$ to $r_1$ the \Info{spread} rises 0.0167 % source: base_spread_weak
$\to$ 0.667 % source: base_spread_strong
while profit at the prior (before any news) falls 4.80 % source: base_profit_prior_weak
$\to$ 4.29: % source: base_profit_prior_strong
deterrence at every fixed belief, and more for the stock to reveal.}
\end{frame}

% ======================================================================
% Frame 7. Trading pays only against a strong incumbent
% ======================================================================
\begin{frame}[label=F7]{Trading pays only against a strong incumbent}
\framesubtitle{Orders $(q_H,q_L)\in[-1,1]$ after a high / low challenger value; trading cost $k$ per unit; prices anticipate entry}
\hypertarget{main:bound}{}\hypertarget{main:orders}{}\hypertarget{main:notation}{}
\vspace{4pt}
{\small The price reveals the market's belief $\mu=\Pr(H)$, kept by noise within $[m,M]$: $m=1/(1+e^{2/b})\approx0.27$, % source: base_m
$M=1-m$ % source: base_M
(Laplace noise, scale $b=2$)\par} % source: base_b
\vspace{3pt}
{\small \textbf{Residual advantage} = what the investor knows a share pays minus its price; entry $\ge\rho$ because cheap preparation pays after any price (the low-cost floor, next frame).\par}
\vspace{2pt}
{\centering\small
\begin{tabular}{@{}l>{\centering\arraybackslash}p{2.5cm}c>{\centering\arraybackslash}p{2.5cm}c>{\centering\arraybackslash}p{2.8cm}@{}}
\toprule
Residual advantage per unit $=$ & entry probability\newline $(\ge\rho)$ & $\times$ & \Info{target-payoff spread $\Delta_T$} & $\times$ & market's remaining uncertainty $(\ge m)$\\
\bottomrule
\end{tabular}\par\vspace{3pt}
$\rho\,m\,\Info{\Delta_T}\ \le\ \text{advantage}\ \le\ \Info{\Delta_T}$\par}
{\small
\begin{itemize}\setlength{\itemsep}{0pt}
\item \textbf{Weak $r_0=1.2$:} % source: base_r_weak
advantage $\le\Delta_T(r_0)=0.0167$% source: base_spread_weak
${}<k=0.02$, % source: base_k
so every order loses; no trade $(0,0)$
\item \textbf{Strong $r_1=3$:} % source: base_r_strong
each extra unit earns $\ge(1-1/b)\,\rho m\Delta_T(r_1)=0.0224>k$, % source: base_k + base_margin_strong_trade (0.02 + 0.00241 = 0.0224)
so full orders $(1,-1)$
\end{itemize}}
% Spoken (plan section 8, frame 7 render fallback): (1 - 1/b) is a haircut for the order's own price impact:
% with Laplace noise, one more unit erodes the per-unit advantage by at most a fraction 1/b.
% source: base_b (2), base_rho (0.25), base_m (m), base_spread_strong (0.667), base_k (0.02), base_margin_strong_trade (0.00241)
\GrayLine{\Info{\textbf{Trading-cost window}} $\Delta_T(r_0)<k<(1-1/b)\,\rho m\Delta_T(r_1)$ (analytical): with $\rho=0.25$, $(1-1/2)\times0.25\times m\times0.667=0.0224=k+{}$theorem margin 0.00241.}
\TakeawayWithNav[0.2em]{The same shift that lowers the challenger's profit switches informed trading on.}
\PlaceNav{\hyperlink{app:bound}{\beamergotobutton{Global trading bound}}\NavSep\hyperlink{app:orders}{\beamergotobutton{Orders, units, costs}}\NavSep\hyperlink{app:notation}{\beamergotobutton{Notation}}}
\end{frame}

% ======================================================================
% Frame 8. X2: only good news makes expensive preparation pay
% ======================================================================
\begin{frame}[label=F8]{Only good news makes expensive preparation pay}
\framesubtitle{$B_r(\mu)=g_L+\mu(g_H-g_L)$ at the worst, prior and best belief a price induces; $c_L=1$ % source: base_c_low
(prob.\ $\rho=0.25$), % source: base_rho
$c_H=6$} % source: base_c_high
\hypertarget{main:suff}{}\hypertarget{main:laplace}{}\hypertarget{main:floor}{}
\begin{columns}[T,onlytextwidth]
\column{0.6\textwidth}
% X2 is drawn at its final size (3.6 x 2.5 in, 10 pt text): included at natural width, no rescaling
\includegraphics[width=0.94\linewidth]{profit_thresholds_talk.pdf}
\column{0.38\textwidth}
\raggedright\small
\vspace{2pt}
\Cost{\textbf{Low-cost floor}}\enspace $c_L<B_{r_1}(m)$: a cheap challenger enters after any price, so entry $\ge\rho$ (the bound used on the previous frame)\par
\vspace{6pt}
\Cost{\textbf{High-cost window}}\enspace $B_{r_0}(1/2)=4.80<c_H=6<B_{r_1}(M)$: % source: base_profit_prior_weak (4.80), base_c_high (6)
the expensive challenger stays out at the weak prior and enters after the best news against $r_1$. Beyond a ceiling strength, even the best price falls short ($r_2$)\par
\vspace{4pt}
\Status{analytical (Props.~A.1, A.2, A.4); costs are inputs}\par
\end{columns}
\NavGap
\PlaceNav{\hyperlink{app:suff}{\beamergotobutton{Reading the price}}\NavSep\hyperlink{app:laplace}{\beamergotobutton{Why Laplace noise}}\NavSep\hyperlink{app:floor}{\beamergotobutton{Is the floor doing the work?}}}
\end{frame}

% ======================================================================
% Frame 9. Proposition 2
% ======================================================================
\begin{frame}[label=F9]{Proposition 2: a stronger incumbent can raise entry}
\framesubtitle{Weak $r_0$, strong $r_1$, stronger still $r_2$; all other primitives equal}
\hypertarget{main:proof}{}\hypertarget{main:margins}{}\hypertarget{main:forces}{}
\begin{ResultBox}[Proposition 2 (analytical), under the three conditions below]
\small
\PropItem{(i)}\textbf{Weak $r_0$:} unique outcome no trade, $(q_H,q_L)=(0,0)$; the price is uninformative; entry $=\rho$
\PropItem{(ii)}\textbf{Strong $r_1$:} unique outcome full orders $(1,-1)$; the price is informative; entry $>\rho$; a high-value challenger acquires the target more often
\PropItem{(iii)}\textbf{Stronger still, $r_2$,} where the low-cost floor and profitable trading persist but $B_{r_2}(M)<c_H$: full orders, entry back to $\rho$\par
\end{ResultBox}
% Conditions (left) and the scope and uniqueness footnote (right) side by side, so the condition rows
% can breathe without pushing the chain into the button band.
\noindent\begin{minipage}[t]{0.585\textwidth}\small
\renewcommand{\arraystretch}{1.15}
\begin{tabular}[t]{@{}l@{\hspace{0.8em}}l@{}}
\Cost{\textbf{Low-cost floor}} & $c_L<B_{r_1}(m)$\\
\Cost{\textbf{High-cost window}} & $B_{r_0}(1/2)<c_H<B_{r_1}(M)$\\
\Info{\textbf{Trading-cost window}} & $\Delta_T(r_0)<k<(1-1/b)\,\rho m\Delta_T(r_1)$
\end{tabular}
\end{minipage}\hfill
\begin{minipage}[t]{0.40\textwidth}\raggedright
{\color{cGrayText}\footnotesize (i)--(ii) on a nonempty open set of primitives, every $h>\ell$; (iii) at the benchmark and nearby. Uniqueness: arbitrary mixed orders; every unilateral deviation $q\in[-1,1]$.\par}
\end{minipage}\par
% Plan render check: the chain wraps beside three buttons, so its first arrow ("stronger incumbent ->") is dropped.
\TakeawayWithNav[0.3em]{\Info{wider $\Delta_T$ $\to$ informed trading pays $\to$ informative price $\to$ good news clears $c_H$ $\to$ the expensive challenger enters}}
\PlaceNav{\hyperlink{app:proof}{\beamergotobutton{Proof logic}}\NavSep\hyperlink{app:margins}{\beamergotobutton{Margins}}\NavSep\hyperlink{app:forces}{\beamergotobutton{Two forces side by side}}}
\end{frame}

% ======================================================================
% Frame 10a. Benchmark table (build step 1)
% ======================================================================
\begin{frame}[label=F10a]{At the benchmark, entry rises from 0.250 to 0.523}
% source (title): base_entry_weak (0.250), base_entry_strong (0.523)
\framesubtitle{Declared benchmark; units arbitrary (inputs in backup); unique equilibrium outcomes (analytical)}
{\footnotesize Entry = Pr(challenger prepares); entry $\rho=0.25$ = only the cheap challenger enters; % source: base_rho
high-value challenger ownership = Pr(high-value challenger acquires the target)\par}
\vspace{4pt}
\makebox[\textwidth][l]{\hspace*{0.7cm}\small
\begin{tabular}{@{}p{6.1cm}>{\centering\arraybackslash}p{2.2cm}>{\centering\arraybackslash}p{2.2cm}@{}}
\toprule
 & weak\newline $r_0=1.2$ % source: base_r_weak
 & \textbf{strong}\newline $\boldsymbol{r_1=3}$\\ % source: base_r_strong
\midrule
Challenger's profit at the prior $B_r(1/2)$ & 4.80 % source: base_profit_prior_weak
 & \textbf{4.29}\\ % source: base_profit_prior_strong
\Info{Target-payoff spread $\Delta_T$} & 0.0167 % source: base_spread_weak
 & \textbf{0.667}\\ % source: base_spread_strong
Investor orders $(q_H,q_L)$ & $(0,0)$ & $\boldsymbol{(1,-1)}$\\
Price & uninformative & \textbf{informative}\\
Entry & 0.250 % source: base_entry_weak
 & \textbf{0.523}\\ % source: base_entry_strong
High-value ownership & 0.125 % source: base_ownership_weak
 & \textbf{0.324}\\ % source: base_ownership_strong
\bottomrule
\end{tabular}}
\Takeaway{Entry rises 27.28 percentage points % source: base_entry_change_pp
while the challenger's profit at the prior falls: the price, not the prize, brings it in.}
\end{frame}

% ======================================================================
% Frame 10b. Benchmark table with the r2 column (build step 2; same frame number)
% ======================================================================
\begin{frame}[noframenumbering,label=F10b]{At the benchmark, entry rises from 0.250 to 0.523}
% source (title): base_entry_weak (0.250), base_entry_strong (0.523)
\framesubtitle{Declared benchmark; units arbitrary (inputs in backup); unique equilibrium outcomes (analytical)}
\hypertarget{main:scale}{}\hypertarget{main:entry}{}\hypertarget{main:cert}{}
{\footnotesize Entry = Pr(challenger prepares); entry $\rho=0.25$ = only the cheap challenger enters; % source: base_rho
high-value challenger ownership = Pr(high-value challenger acquires the target)\par}
\vspace{4pt}
\makebox[\textwidth][l]{\hspace*{0.7cm}\small
\begin{tabular}{@{}p{6.1cm}>{\centering\arraybackslash}p{2.2cm}>{\centering\arraybackslash}p{2.2cm}>{\centering\arraybackslash}p{2.2cm}@{}}
\toprule
 & weak\newline $r_0=1.2$ % source: base_r_weak
 & \textbf{strong}\newline $\boldsymbol{r_1=3}$ % source: base_r_strong
 & stronger still\newline $r_2=3.6$\\ % source: base_r_collapse
\midrule
Challenger's profit at the prior $B_r(1/2)$ & 4.80 % source: base_profit_prior_weak
 & \textbf{4.29} % source: base_profit_prior_strong
 & \textit{\textcolor{cGrayText}{lower still}}\\
\Info{Target-payoff spread $\Delta_T$} & 0.0167 % source: base_spread_weak
 & \textbf{0.667} % source: base_spread_strong
 & \textit{\textcolor{cGrayText}{wider still}}\\
Investor orders $(q_H,q_L)$ & $(0,0)$ & $\boldsymbol{(1,-1)}$ & $(1,-1)$\\
Price & uninformative & \textbf{informative} & informative\\
Entry & 0.250 % source: base_entry_weak
 & \textbf{0.523} % source: base_entry_strong
 & 0.250\\ % source: base_entry_collapse
High-value ownership & 0.125 % source: base_ownership_weak
 & \textbf{0.324} % source: base_ownership_strong
 & 0.125\\ % source: tables/table2_equilibrium_controls.tex Panel A
\bottomrule
\end{tabular}}
\TakeawayWithNav{Rise, then fall: past a ceiling strength even the best price cannot cover $c_H$.}
\PlaceNav{\hyperlink{app:scale}{\beamergotobutton{Benchmark scale}}\NavSep\hyperlink{app:entry}{\beamergotobutton{How entry is computed}}\NavSep\hyperlink{app:cert}{\beamergotobutton{Which equilibrium?}}}
\end{frame}

% ======================================================================
% Frame 11. Freeze the information
% ======================================================================
\begin{frame}[label=F11]{Freeze the information and deterrence returns}
\framesubtitle{Entry at $r_0=1.2$ % source: base_r_weak
and $r_1=3$; % source: base_r_strong
equilibria of the feedback game, then information controls}
\hypertarget{main:controls}{}\hypertarget{main:welfare}{}
\begin{center}\small
\begin{tabular}{@{}m{5.9cm}ccc>{\raggedright\arraybackslash}m{3.9cm}@{}}
\toprule
 & $r_0=1.2$ & $r_1=3$ & Direction & Status\\ % source: base_r_weak, base_r_strong
\midrule
Equilibrium of the feedback game\newline{\footnotesize (price seen, trading re-solved)} & 0.250 % source: base_entry_weak
 & \textbf{0.523} % source: base_entry_strong
 & $\uparrow$ & \StatusF{analytical}\\
\midrule
\multicolumn{5}{@{}l}{\emph{Information controls}}\\
Frozen informative orders $(1,-1)$\newline{\footnotesize (imposed at both strengths; control, not an equilibrium at $r_0$)} & 0.562 % source: base_frozen_entry_weak
 & 0.523 % source: base_frozen_entry_strong
 & $\downarrow$ & \StatusF{numerical diagnostic;\newline sign analytical (Prop.~A.3)}\\
Price hidden from challenger\newline{\footnotesize (equilibrium of the no-price-access game)} & 0.250 % source: base_hidden_entry_weak
 & 0.250 % source: base_hidden_entry_strong
 & flat & \StatusF{analytical}\\
\bottomrule
\end{tabular}
\end{center}
\vspace{-2pt}
\GrayLine{The price-hidden row is implied by the high-cost window. Without the price, expensive preparation never pays and only the entry floor remains.}
\TakeawayWithNav{Hold the price's information fixed and a stronger incumbent weakly lowers entry.}
\PlaceNav{\hyperlink{app:controls}{\beamergotobutton{Controls in detail}}\NavSep\hyperlink{app:welfare}{\beamergotobutton{Price level or information?}}}
\end{frame}

% ======================================================================
% Frame 12. Coexistence
% ======================================================================
\begin{frame}[label=F12]{At intermediate strength, equilibria coexist}
\framesubtitle{Computer-assisted (Proposition 3): interval arithmetic on exact decimal inputs}
\hypertarget{main:corr}{}
\vspace{4pt}
\begin{center}\small
\renewcommand{\arraystretch}{1.2}
\begin{tabular*}{0.92\textwidth}{@{\extracolsep{\fill}}cccc@{}}
\toprule
Strength $r$ & $q_H$ & $q_L$ (certified) & Entry (certified)\\
\midrule
1.55 % source: cert_a_r
 & 1 & $\approx-0.460$ % source: cert_a_v_interval (negated)
 & $\approx0.545$\\ % source: cert_a_entry_interval
1.60 % source: cert_b_r
 & 1 & $\approx-0.707$ % source: cert_b_v_interval (negated)
 & $\approx0.549$\\ % source: cert_b_entry_interval
1.65 % source: cert_c_r
 & 1 & $\approx-0.903$ % source: cert_c_v_interval (negated)
 & $\approx0.551$\\ % source: cert_c_entry_interval
\bottomrule
\end{tabular*}
\end{center}
\vspace{-4pt}
\begin{itemize}\setlength{\itemsep}{3pt}
\item Buy fully after good news, sell partially after bad news; certified entry rises from one strength to the next, above the 0.523 at $r_1$ % source: base_entry_strong
\item Each economy also has a no-trade equilibrium with entry $\rho=0.25$ % source: base_rho
\Status{(analytical, Prop.~A.4)}
\item Elsewhere the numerical search is not exhaustive; the full correspondence is open
\end{itemize}
% Spoken: the theorem compares r0 with r1, where trading is unique; it does not say entry rises
% monotonically in between (A24).
\TakeawayWithNav{At three certified strengths an informative equilibrium coexists with no trade; I claim nothing about entry between them.}
\PlaceNav{\hyperlink{app:corr}{\beamergotobutton{Correspondence}}}
\end{frame}

% ======================================================================
% Frame 13. Robustness
% ======================================================================
\begin{frame}[label=F13]{The entry reversal survives four model changes}
\framesubtitle{Entry, weak to strong; each row its own analytical result and declared parameters (Prop.~2 (i)--(ii))}
\hypertarget{main:robust}{}
\vspace{2pt}
\begin{center}\small
\begin{tabular}{@{}>{\raggedright\arraybackslash}p{10.4cm}cc@{}}
\toprule
Specification & Weak & Strong\\
\midrule
Benchmark: Laplace noise, two cost levels & 0.250 % source: base_entry_weak
 & 0.523\\ % source: base_entry_strong
Atomless costs, half-width 0.1 % source: cost_halfwidth
 & 0.250 % source: tables/table3_extensions.tex Panel A, "Laplace, cost mixture"
 & 0.523\\ % source: cost_mix_laplace_entry_strong
Logistic noise & 0.250 % source: tables/table3_extensions.tex Panel A, "logistic, cost atoms"
 & 0.302\\ % source: logistic_entry_strong
Small value gap, $h=2$, % source: moderate_h
$\ell=1$ % source: moderate_ell
 & 0.250 % source: moderate_entry_weak
 & 0.527\\ % source: moderate_entry_strong
Complementary signals (declared example): investor accuracy 0.70, % source: signal_trader_accuracy_value
challenger accuracy 0.75; % source: signal_buyer_accuracy_value
$\rho=0.85$, % source: signal_rho
strengths 1.1 % source: signal_r_weak
vs 2.3 % source: signal_r_strong
 & 0.850 % source: signal_entry_weak
 & 0.879\\ % source: signal_entry_strong
\bottomrule
\end{tabular}
\end{center}
\vspace{-3pt}
% source: tables/table_signal_grid.tex (25-cell grid; 6 met, 9 unchanged with d <= 0.75, 10 fall with d >= 0.76)
{\color{cGrayText}\footnotesize Rows differ in $\rho$ and strengths; compare within a row, not across rows.\par
\vspace{2pt}
Signal grid, 25 accuracy cells: entry rises in the 6 meeting Prop.~A.7 (analytical) and is unchanged in 9 with challenger accuracy ${}\le0.75$; it falls about 4~pp when challenger accuracy is ${}\ge0.76$, because its own good signal already triggers entry against the weak incumbent (numerical diagnostic).\par}
\TakeawayWithNav[0.5em]{The $r_0\to r_1$ reversal survives atomless costs, logistic noise, a small value gap and complementary private signals (declared example).}
\PlaceNav{\hyperlink{app:robust}{\beamergotobutton{Full table}}}
\end{frame}

% ======================================================================
% Frame 14. Price access and welfare
% ======================================================================
\begin{frame}[t,label=F14]{Price access raises proceeds and surplus at $r_1$}
\framesubtitle{Incumbent held at $r_1=3$; % source: base_r_strong
price seen or hidden; trading re-solved in both (Prop.~A.9, analytical)}
\hypertarget{main:reserve}{}
\vspace{6pt}
\begin{center}
\renewcommand{\arraystretch}{1.25}
\begin{tabular*}{0.92\textwidth}{@{\extracolsep{\fill}}lccc@{}}
\toprule
 & Price observed & Price hidden & Gain\\
\midrule
Expected target proceeds & 0.872 % source: base_revenue_feedback
 & 0.615 % source: base_revenue_hidden
 & 0.258\\ % source: base_revenue_gain
Acquisition surplus net of preparation costs & 2.38 % source: tables/table2_equilibrium_controls.tex Panel C
 & 2.30 % source: tables/table2_equilibrium_controls.tex Panel C
 & 0.0802\\ % source: base_net_surplus_gain
\bottomrule
\end{tabular*}
\end{center}
\begin{itemize}\setlength{\itemsep}{4pt}
\item Same full orders in both economies, so trading costs coincide
\item Each extra entry happens only when its expected profit covers its cost
\end{itemize}
% source: base_revenue_feedback (0.872392), base_revenue_hidden (0.614583), base_revenue_gain (0.257809): 0.872392 - 0.614583 = 0.257809
\vskip 0pt plus 4filll% outweighs the frame's own bottom fill, so the note sits just above the button band
\GrayLine{Gain from unrounded values. At $r_0$ prices are uninformative, so there is no gain. The comparison holds the sale rule and the incumbent distribution fixed.}
\NavGap
\PlaceNav{\hyperlink{app:reserve}{\beamergotobutton{Reserve}}}
\end{frame}

% ======================================================================
% Frame 15. Three predictions (paper Section 7)
% ======================================================================
\begin{frame}[label=F15]{Three predictions}
\framesubtitle{Public sale processes with one prepared bidder and a further buyer still deciding}
\begin{itemize}\setlength{\itemsep}{8pt}
\item \textbf{A stronger bidder can draw a buyer in} when the price comes before the buyer's decision; without the price, a stronger bidder weakly lowers entry \Status{(analytical, Prop.~2; Prop.~A.3)}
\item \textbf{Buyers that prepare after good news win more often:} among buyers that prepare at a price, the share of high-value buyers is the price-implied posterior
\item \textbf{Against a very strong bidder, entry stops responding} to the price, while the price stays informative \Status{(analytical, Prop.~2(iii))}
\end{itemize}
\medskip
\GrayLine{Outcome: the start of diligence, not public offers. Pre-bid returns also reflect anticipation of the bid, so each test dates the price before the decision (Schwert 1996; Online Appendix D).}
\end{frame}

% ======================================================================
% Frame 16. Conclusion (no navigation)
% ======================================================================
\begin{frame}[label=F16]{Conclusion}
\begin{itemize}\setlength{\itemsep}{20pt}
\item \KeyIdea{A stronger incumbent can bring the challenger in:} at the benchmark, entry rises from 0.250 % source: base_entry_weak
to 0.523 % source: base_entry_strong
although the challenger's expected profit before any news falls
\item \textbf{Why:} competition makes target proceeds more sensitive to who the challenger is; informed trading puts that into the price; good news draws in the expensive challenger. Freeze that information and deterrence returns.
\item \textbf{Implication:} when the target trades, a stronger rival is not a pure deterrent, because who competes depends on what the price reveals before anyone prepares.
\end{itemize}
\end{frame}

% ######################################################################
% APPENDIX
% ######################################################################
\AppendixStart

% ---------------------------------------------------------------- A1
\begin{frame}[label=A1]{Which deals have a decision interval}
\hypertarget{app:interval}{}
\begin{itemize}\setlength{\itemsep}{6pt}
\item A disclosed approach, an announced strategic review or an open contest can create one if the opportunity becomes public before the buyer set is fixed
\item Only then can a prospective challenger watch the price while it decides
\item Strength = the public distribution of the incumbent's value, not an announced bid
\item Imprivata's proxy separates an unsolicited approach, outreach to potential buyers, and indications of interest conditional on diligence, which documents a costly preparation stage; the three predictions say how to test whether prices draw buyers into it \graycite{Imprivata 2016; Boone and Mulherin 2007; Gentry and Stroup 2019}
\end{itemize}
\vspace{0.4em}
\GrayLine{Whether a deal fits is a question about its chronology; Online Appendix D shows how to date it (Section 2.1).}
\BackButton{main:interval}
\end{frame}

% ---------------------------------------------------------------- A2
\begin{frame}[label=A2]{Evidence: the pilot design}
\hypertarget{app:evidence}{}
\begin{itemize}\setlength{\itemsep}{6pt}
\item \textbf{Pilot (Online Appendix D):} reconstruct from disclosure records the timing of approaches, public visibility, buyer contacts, diligence, proposals and final selection
\item Separate when an event occurred from when it first became public
\item \textbf{First requirement:} a publicly understood sale opportunity that stays contestable while a challenger decides
\item \textbf{Outcome:} the decision to prepare; prices must precede it
\item \textbf{Caution:} returns before a bid also reflect its anticipation \graycite{Schwert 1996}
\end{itemize}
\vspace{0.4em}
\GrayLine{Online Appendix D sets out the pilot design for the predictions of Section 7.}
\BackButton{main:evidence}
\end{frame}

% ---------------------------------------------------------------- A3
\begin{frame}[label=A3]{Related work in more detail}
\framesubtitle{Increment: under one sale rule the two returns move oppositely, which produces the entry reversal}
\hypertarget{app:lit}{}
\vspace{3pt}
{\scriptsize
\setlength{\tabcolsep}{3pt}
\renewcommand{\arraystretch}{1.0}
\begin{tabular}{@{}>{\raggedright\arraybackslash}p{5.1cm}>{\raggedright\arraybackslash}p{4.8cm}>{\raggedright\arraybackslash}p{4.7cm}@{}}
\toprule
Work & Shows & Here\\
\midrule
Dow, Goldstein and Guembel (2017) & investment feeds back on information & one surplus, two opposed claims\\\addlinespace[1pt]
Edmans, Goldstein and Jiang (2015) & real actions curb bad-news trading & rivals shift claim sensitivity\\\addlinespace[1pt]
Edmans, Goldstein and Jiang (2012) & prices affect takeover activity & a challenger learns its own value\\\addlinespace[1pt]
Fishman (1988);\newline Hirshleifer and Png (1989) & stronger rivals deter preparation & kept at fixed information (Prop.~A.3)\\\addlinespace[1pt]
Levin and Smith (1994);\newline Gentry and Stroup (2019) & the bidder pool is endogenous & entry responds to the price\\\addlinespace[1pt]
Roberts and Sweeting (2013) & selective entry shapes procedure & sale terms left open\\\addlinespace[1pt]
Persico (2000) & formats shape bidders' information & the informed party is not a bidder\\\addlinespace[1pt]
Luo (2005) & learning from announcement returns & learning precedes entry\\\addlinespace[1pt]
Betton et al.\ (2014); Lin et al.\ (2025) & negotiation feedback; payment choice & cash fixed; entry of a new bidder\\\addlinespace[1pt]
Cornelli and Li (2002) & arbitrage positions affect tenders & the investor does not tender\\\addlinespace[1pt]
Pernoud and Gleyze (2026);\newline Liu and Bernhardt (2022);\newline Carlin et al.\ (2026) & learning own values and rivals;\newline post-auction feedback;\newline bidder-pool choice & the informed party trades outside the auction, before entry\\
\bottomrule
\end{tabular}\par}
\BackButton{main:lit}
\end{frame}

% ---------------------------------------------------------------- A4
\begin{frame}[label=A4]{What the model leaves out, and why}
\hypertarget{app:omit}{}
\vspace{4pt}
\begin{center}\small
\renewcommand{\arraystretch}{1.25}
\begin{tabular}{@{}>{\raggedright\arraybackslash}p{2.6cm}>{\raggedright\arraybackslash}p{11.0cm}@{}}
\toprule
\textbf{In the model} & Neither bidder trades. The investor has no initial position, cannot bid, tender or acquire, and has no control rights. The sale binds all shares (no tendering or holdout). Every wrong-signed order loses against every candidate schedule, so the investor cannot profit by faking good news.\\
\midrule
\textbf{Outside the model} & Toeholds \graycite{Bulow, Huang and Klemperer 1999}; free riding \graycite{Grossman and Hart 1980}; announced or jump bids as signals \graycite{Fishman 1988} (strength here is a distribution); incumbent or target manipulation and endogenous investor research \graycite{Goldstein and Guembel 2008}; first-price and bargaining-with-trading equilibria.\\
\bottomrule
\end{tabular}
\end{center}
\Takeaway{The paper does not solve toeholds, so it makes no statement on the direction of a toehold effect.}
\BackButton{main:omit}
\end{frame}

% ---------------------------------------------------------------- A5
\begin{frame}[label=A5]{Complementary, not superior, information}
\hypertarget{app:info}{}
\begin{itemize}\setlength{\itemsep}{6pt}
\item The challenger knows its own integration plans; a specialist investor may know the target's customers, technology and product demand
\item Different facts about the same acquisition match
\item Preparation (diligence) is what reveals the exact value to the challenger
\item The benchmark's perfectly informed investor and uninformed challenger make the mechanism visible; Section 5.3 removes both extremes
\item The reversal holds with challenger accuracy $d=0.75$ % source: signal_buyer_accuracy_value
above investor accuracy $a=0.70$ % source: signal_trader_accuracy_value
\Status{(Prop.~A.7, analytical)}
\end{itemize}
\NavGap
\PlaceNav{\hyperlink{app:signals}{\beamergotobutton{The signal economy}}\NavSep\hyperlink{main:info}{\beamerreturnbutton{Back}}}
\end{frame}

% ---------------------------------------------------------------- A6
\begin{frame}[label=A6]{The price helps a challenger with its own signal}
\hypertarget{app:signals}{}
{\small
\begin{itemize}\setlength{\itemsep}{3pt}
\item Investor signal $T$ (accuracy $a$), challenger signal $Y$ (accuracy $d$); the challenger uses $\Pr(H\mid P,\,Y=y)$, $P$ = stock price
\item Entry is state dependent: $e_H$, $e_L$ = entry probability when the challenger is high / low value; the price still reveals the market posterior
\item Declared example $a=0.70$, % source: signal_trader_accuracy_value
$d=0.75$: % source: signal_buyer_accuracy_value
entry 0.850 % source: signal_entry_weak
$\to$ 0.879 % source: signal_entry_strong
\Status{(analytical, Prop.~A.7)}; $\rho=0.85$, % source: signal_rho
$c_H=7.14$, % source: signal_c_high
$k=0.015$, % source: signal_k
strengths 1.1 % source: signal_r_weak
vs 2.3 % source: signal_r_strong
\end{itemize}}
\vspace{2pt}
% source: tables/table_signal_grid.tex (all 25 cells; counts and pp ranges as printed)
\begin{center}\small
\begin{tabular}{@{}clcl@{}}
\toprule
Cells & Accuracy grid (25 cells) & Entry change (pp) & Status\\
\midrule
6 & meet Prop.~A.7 & rises 2.81 to 3.09 & \Status{analytical}\\
9 & $d\le0.75$, condition not met & unchanged & \Status{numerical diagnostic}\\
10 & $d\ge0.76$ & falls 4.03 to 4.49 & \Status{numerical diagnostic}\\
\bottomrule
\end{tabular}
\end{center}
\Takeaway{The price matters when the challenger's own signal is informative but not decisive.}
\BackButton{app:info}
\end{frame}

% ---------------------------------------------------------------- A7
\begin{frame}[label=A7]{Equilibrium and the two comparisons}
\hypertarget{app:eq}{}
\begin{itemize}\setlength{\itemsep}{6pt}
\item \textbf{Definition:} order distributions $\sigma_H,\sigma_L$ (mixing allowed; any deviation $q\in[-1,1]$); a rational price that anticipates entry; Bayesian beliefs from the price; optimal entry; truthful bids
\item Order flow $X=q+Z$, Laplace noise $Z$ with scale $b$; stock price $P(X)$ = expected target payoff given $X$
\item \textbf{Two comparisons kept apart:} a unilateral deviation is evaluated against fixed price and entry schedules; a cross-economy comparison re-solves them
\item \textbf{Uniqueness} refers to trading and on-path entry under truthful bidding
\end{itemize}
\BackButton{main:eq}
\end{frame}

% ---------------------------------------------------------------- A8
\begin{frame}[label=A8]{Acquisition payoffs in closed form}
\hypertarget{app:payoffs}{}
{\small
\abovedisplayskip=0pt\belowdisplayskip=4pt
\begin{align*}
&t_0=p(1-p/r), && t_H=r/2+p^2/(2r), && t_L=\ell-(\ell^2-p^2)/(2r),\\
&g_H=h-r/2-p^2/(2r), && g_L=(\ell^2-p^2)/(2r), && \Info{\Delta_T=(r-\ell)^2/(2r)}
\end{align*}\par}
% source: tables/table1_auction_primitives.tex; 3 significant digits
{\centering\footnotesize
\renewcommand{\arraystretch}{0.95}
\begin{tabular}{@{}lcc@{}}
\toprule
 & weak $r_0=1.2$ & strong $r_1=3$\\ % source: base_r_weak, base_r_strong
\midrule
Proceeds without a challenger, $t_0$ & 0.292 & 0.417\\ % source: tables/table1_auction_primitives.tex
Proceeds with a high-value challenger, $t_H$ & 0.704 & 1.54\\ % source: tables/table1_auction_primitives.tex
Proceeds with a low-value challenger, $t_L$ & 0.688 & 0.875\\ % source: tables/table1_auction_primitives.tex
Gross profit of a high-value challenger, $g_H$ & 9.30 & 8.46\\ % source: tables/table1_auction_primitives.tex
Gross profit of a low-value challenger, $g_L$ & 0.313 & 0.125\\ % source: tables/table1_auction_primitives.tex
\Info{Target-payoff spread, $\Delta_T=t_H-t_L$} & 0.0167 & 0.667\\ % source: base_spread_weak, base_spread_strong
Expected gross profit at the prior, $B_r(1/2)$ & 4.80 & 4.29\\ % source: base_profit_prior_weak, base_profit_prior_strong
\bottomrule
\end{tabular}\par}
\vspace{4pt}
\GrayLine{Proposition 1 (analytical) holds for any $F$ on $[0,\bar r]$, $\ell<\bar r<h$; strict when the change in $F$ has positive integral over the relevant range.}
\Takeaway[0.4em]{The closed forms are the uniform case; the two signs are not a uniform-distribution artifact.}
\BackButton{main:payoffs}
\end{frame}

% ---------------------------------------------------------------- A9
\begin{frame}[label=A9]{The payment rule decides the sign of the spread effect}
\hypertarget{app:payrule}{}
\begin{columns}[T,onlytextwidth]
\column{0.53\textwidth}
% X4 is drawn at its final size (3.3 x 2.3 in, 10 pt text): included at natural width
\includegraphics[width=0.94\linewidth]{bargaining_spread_talk.pdf}
\column{0.45\textwidth}
\raggedright\small
\vspace{4pt}
Zero-reserve verifiable-value institution, Nash weight $\eta$: seller gets $t_\eta=(1-\eta)\cdot{}$runner-up value${}+\eta\cdot{}$winner value\par
\vspace{6pt}
Stronger incumbent: challenger profit falls weakly for every $\eta$. Spread $\Info{\Delta_\eta}=\eta(h-\ell)+(1-2\eta)\,\mathbb{E}[(R-\ell)_+]$ rises (weakly) if $\eta<1/2$, falls (weakly) if $\eta>1/2$, for any $R\sim F$ on $[0,\bar r]$, $\ell<\bar r<h$, $0\le\eta<1$ \Status{(Prop.~A.8, analytical)}\par
\end{columns}
\Takeaway{Payment stage only: an entry reversal under bargaining is not solved, and first-price equilibria are not solved.}
\BackButton{main:payrule}
\end{frame}

% ---------------------------------------------------------------- A10
\begin{frame}[label=A10]{Why trading is unique: a global bound}
\hypertarget{app:bound}{}
\begin{itemize}\setlength{\itemsep}{4pt}
\item \textbf{Weak:} gross advantage per unit $\le\Delta_T(r_0)<k$, so every nonzero order loses and the posterior stays at $1/2$
\item \textbf{Strong:} with $\Pi_\theta(s)$ the expected residual per unit at order size $s$,
\end{itemize}
\vspace{-4pt}
\[
U_\theta(s)=s\,\Pi_\theta(s)-ks,\qquad U_\theta'(s)\ \ge\ (1-1/b)\,\rho m\,\Info{\Delta_T(r_1)}-k\ >\ 0\quad\text{on }[0,1]
\]
\vspace{-8pt}
\begin{itemize}\setlength{\itemsep}{4pt}
\item Holds against every candidate price and entry schedule, including mixed orders: global, not first-order
\item \textbf{Investor manipulation (in-model):} both residual advantages in eq.~(11) are positive under every candidate order distribution, so a wrong-signed order, such as a low-value investor buying to fake good news and trigger entry, has negative gross payoff and still pays $k$
\item Incumbent or target manipulation and toeholds: outside the model (A4)
\end{itemize}
\BackButton{main:bound}
\end{frame}

% ---------------------------------------------------------------- A11
\begin{frame}[label=A11]{Orders, units and the trading cost}
\hypertarget{app:orders}{}
\begin{itemize}\setlength{\itemsep}{6pt}
\item $q\in[-1,1]$ is a normalized small trading unit, not ownership of the target, measured in the same units as the noise (scale $b$); the investor has no initial position
\item $k$ is an execution or position-carrying friction, distinct from adverse-selection price impact, which competitive pricing already generates
\item \textbf{Why a corner, unlike Kyle:} the per-unit advantage is at least $\rho m\Delta_T$ because beliefs stay in $[m,M]$, and one more unit erodes it by at most a fraction $1/b$ (the haircut)
\item So marginal profit stays above $k$ over the whole interval, and full orders are the unique best response (A10)
\item Against the weak incumbent the same bounds make every order lose
\end{itemize}
\BackButton{main:orders}
\end{frame}

% ---------------------------------------------------------------- A12
\begin{frame}[label=A12]{Notation}
\hypertarget{app:notation}{}
\vspace{4pt}
% Rows in the order the symbols are introduced in the talk (structure plan, section 6.1)
{\footnotesize
\begin{columns}[T]
\column{0.49\textwidth}
\begin{tabular}{@{}l>{\raggedright\arraybackslash}p{4.6cm}@{}}
\toprule
Symbol & Meaning\\
\midrule
$R$; $r$ & incumbent value $R\sim U[0,r]$; $r$ = incumbent strength\\
$\theta\in\{H,L\}$; $h,\ell$ & challenger quality; worth $h$ or $\ell$\\
$p$ & reserve price (not the stock price)\\
$C$; $c_L,c_H$; $\rho$ & preparation cost; cheap, expensive; $\Pr(C=c_L)$\\
$t_H,t_L$ & expected target proceeds when a high / low-value challenger enters\\
$\Delta_T$ & target-payoff spread $t_H-t_L$\\
$g_H,g_L$ & gross acquisition profit of a high / low-value challenger\\
\bottomrule
\end{tabular}
\column{0.49\textwidth}
\begin{tabular}{@{}l>{\raggedright\arraybackslash}p{4.6cm}@{}}
\toprule
Symbol & Meaning\\
\midrule
$r_0,r_1$ & weak, strong incumbent\\
$\mu$ & market's belief $\Pr(H)$ revealed by the price\\
$m,M$ & bounds on $\mu$: $m=1/(1+e^{2/b})$, $M=1-m$\\
$b$ & scale of the Laplace noise\\
$k$ & trading cost per unit\\
$(q_H,q_L)$ & orders after a high / low challenger value\\
$B_r(\mu)$ & challenger's expected gross profit at belief $\mu$\\
$r_2$ & stronger still incumbent\\
\bottomrule
\end{tabular}
\end{columns}}
\BackButton{main:notation}
\end{frame}

% ---------------------------------------------------------------- A13
\begin{frame}[label=A13]{The challenger can read the market's belief off the price}
\hypertarget{app:suff}{}
\begin{itemize}\setlength{\itemsep}{3pt}
\item Under any mixed orders the posterior from order flow is bounded (Laplace likelihood ratio within $e^{\pm2/b}$):
\end{itemize}
\vspace{-6pt}
\begin{gather*}
\mu_X(x)=\Pr(H\mid X=x)\in[m,M]\\
\text{price}=t_0+(\text{entry probability at that price})\times(t_L-t_0+\Delta_T\,\mu)
\end{gather*}
\vspace{-6pt}
\begin{itemize}\setlength{\itemsep}{3pt}
\item $t_0$ = expected target proceeds without a challenger
\item The price is strictly increasing in $\mu$, because entry $\ge\rho>0$ and $\Delta_T>0$
\item Atoms and entry jumps allowed; no differentiability needed \Status{(Props.~A.1--A.2, analytical)}
\end{itemize}
\Takeaway{The price is a sufficient statistic for the market's belief, so the challenger loses nothing by seeing the price rather than the flow.}
\BackButton{main:suff}
\end{frame}

% ---------------------------------------------------------------- A14
\begin{frame}[label=A14]{Why Laplace noise, and what logistic noise changes}
\hypertarget{app:laplace}{}
\centering
% X3 rebuilt at its final size (5.8 x 1.7 in, 10 pt sans text): included at natural width, no rescaling
\includegraphics[width=5.8in]{posterior_tail_entry_talk.pdf}
\par\vspace{-2pt}\raggedright
{\footnotesize
\begin{itemize}\setlength{\itemsep}{0pt}
\item Both noise densities $f$ satisfy $|f'|\le f/b$, so the global trading bounds hold for both
\item Laplace: the posterior reaches $[m,M]$ at finite flow; logistic: only as flow $\to\infty$ ($x^*\approx5.42$, % source: logistic_flow_threshold
about 1.5 noise s.d.\ from the center) % source: logistic_threshold_noise_sd
\item Entry 0.302 % source: logistic_entry_strong
vs 0.523; % source: base_entry_strong
atomless costs ($\varepsilon_C=0.1$): % source: cost_halfwidth
0.523 % source: cost_mix_laplace_entry_strong
and 0.301. % source: cost_mix_logistic_entry_strong
Common scale $b$, not common variance; not a Blackwell ranking \Status{(analytical, Props.~A.5--A.6)}
\end{itemize}}
\BackButton{main:laplace}
\end{frame}

% ---------------------------------------------------------------- A15
\begin{frame}[label=A15]{Why the model needs a low-cost floor}
\hypertarget{app:floor}{}
\begin{itemize}\setlength{\itemsep}{6pt}
\item The \Cost{low-cost floor} keeps entry $\ge\rho$ after every price: the price reveals the posterior, the investor's edge is at least $\rho m\Delta_T$, and information can only add entry
\item If cheap preparation fails only at bad prices, bad prices deter it and flows can pool at one price
\item If cheap preparation fails at the prior, no preparation and no trade form an equilibrium
\item Theorem margin at the benchmark: 1.37 % source: base_margin_low_cost
\item The floor alone gives entry 0.250 % source: base_hidden_entry_weak, base_hidden_entry_strong
at both strengths when the price is hidden
\item Atomless cost supports keep it \Status{(Prop.~A.6, analytical)}
\end{itemize}
\BackButton{main:floor}
\end{frame}

% ---------------------------------------------------------------- A16
\begin{frame}[label=A16]{Proof logic in four steps}
\hypertarget{app:proof}{}
\setbeamertemplate{enumerate items}{\textbf{\insertenumlabel.}}
\begin{enumerate}\setlength{\itemsep}{5pt}
\item \textbf{Weak economy:} no order covers $k$; no trade; entry $\rho$
\item \textbf{Strong economy:} the global bound (A10) forces full orders
\item Under full orders the high-cost window puts the belief threshold inside $(1/2,M)$, so good news crosses it with positive probability; $\alpha_H>\alpha_L$ (good news is likelier when the challenger is high value, $\alpha_\theta=\Pr(\text{order flow crosses the entry threshold}\mid\theta)$) tilts extra entry to high values
\item \textbf{Blackwell:} a constant kernel maps the strong experiment to the weak one; no state-independent kernel maps the constant experiment to the strong one
\end{enumerate}
\medskip
\textbf{Part (iii):} with the floor and profitable trading still in place at $r_2$, $B_{r_2}(M)<c_H$ removes expensive entry.
\Takeaway{What does the work: step 2, a bound that holds against every candidate schedule.}
\BackButton{main:proof}
\end{frame}

% ---------------------------------------------------------------- A17
\begin{frame}[label=A17]{Conditions are jointly satisfiable for every $h>\ell$}
\hypertarget{app:margins}{}
\vspace{4pt}
\begin{center}\small
\begin{tabular}{@{}llr@{}}
\toprule
\multicolumn{3}{@{}l}{Theorem margins at the benchmark}\\
\midrule
\Cost{Low-cost floor} & $B_{r_1}(m)-c_L$ & 1.37\\ % source: base_margin_low_cost
\Cost{High-cost window}, weak prior & $c_H-B_{r_0}(1/2)$ & 1.20\\ % source: base_margin_high_prior
\Cost{High-cost window}, best strong price & $B_{r_1}(M)-c_H$ & 0.217\\ % source: base_margin_high_ceiling
\Info{Trading-cost window}, weak & $k-\Delta_T(r_0)$ & 0.00333\\ % source: base_margin_weak_trade
\Info{Trading-cost window}, strong & $(1-1/b)\rho m\Delta_T(r_1)-k$ & 0.00241\\ % source: base_margin_strong_trade
\midrule
Minimum & & 0.00241\\ % source: base_minimum_theorem_margin
\bottomrule
\end{tabular}
\end{center}
\vspace{-4pt}
{\small
\begin{itemize}\setlength{\itemsep}{2pt}
\item Nonemptiness for every $h>\ell$ comes from a construction: take $r_0<r_1$ close to $\ell$, then $k$, $c_H$, $c_L$ strictly inside their windows
\item A mathematical construction, not slack at the benchmark and not an effect-size claim; moderate values $h=2$: % source: moderate_h
minimum margin $8.01\times10^{-4}$ % source: moderate_minimum_theorem_margin
\item Part (iii) is not claimed for every $h>\ell$
\end{itemize}}
\BackButton{main:margins}
\end{frame}

% ---------------------------------------------------------------- A18
\begin{frame}[label=A18]{Two forces from one payment rule}
\hypertarget{app:forces}{}
\vspace{4pt}
\begin{center}\small
\renewcommand{\arraystretch}{1.2}
\begin{tabular}{@{}>{\raggedright\arraybackslash}p{3.3cm}>{\raggedright\arraybackslash}p{4.4cm}>{\raggedright\arraybackslash}p{5.6cm}@{}}
\toprule
 & Deterrence & \Info{Information}\\
\midrule
What a stronger incumbent moves & challenger profit $B_r(\mu)\downarrow$ at every belief & \Info{target-payoff spread $\Delta_T\uparrow$}\\
Who responds & the challenger, at a given belief & the investor, then the price, then the challenger's belief\\
What it needs & nothing & entry after any price (\Cost{low-cost floor}); price seen before entry\\
Effect on entry & \textbf{weakly $\downarrow$} \Status{(Prop.~A.3)} & $\uparrow$ once trading pays and good news clears \Cost{$c_H$}\\
\bottomrule
\end{tabular}
\end{center}
\Takeaway{\Info{stronger incumbent $\to$ wider $\Delta_T$ $\to$ informed trading pays $\to$ informative price $\to$ good news clears} \Cost{$c_H$} \Info{$\to$ the expensive challenger enters}}
\BackButton{main:forces}
\end{frame}

% ---------------------------------------------------------------- A19
\begin{frame}[label=A19]{Benchmark scale: declared, not calibrated}
\hypertarget{app:scale}{}
\vspace{4pt}
{\small
\begin{columns}[c]
\column{0.38\textwidth}
\centering
\begin{tabular}{@{}lcc@{}}
\toprule
Input & Benchmark & Moderate\\
\midrule
$h$ & 10 & 2\\ % source: base_h, moderate_h
$\ell$ & 1 & 1\\ % source: base_ell, moderate_ell
$p$ & 0.5 & 0.5\\ % source: base_p, moderate_p
$\rho$ & 0.25 & 0.25\\ % source: base_rho, moderate_rho
$c_L$ & 1 & 0.3\\ % source: base_c_low, moderate_c_low
$c_H$ & 6 & 0.89\\ % source: base_c_high, moderate_c_high
$b$ & 2 & 2\\ % source: base_b, moderate_b
$k$ & 0.02 & 0.002\\ % source: base_k, moderate_k
\bottomrule
\end{tabular}
\column{0.60\textwidth}
\centering
\begin{tabular}{@{}lcc@{}}
\toprule
 & Benchmark & Moderate\\
\midrule
Strengths & $1.2\to3$ & $1.05\to1.5$\\ % source: base_r_weak, base_r_strong, moderate_r_weak, moderate_r_strong
Entry & $0.250\to0.523$ & $0.250\to0.527$\\ % source: base_entry_weak, base_entry_strong, moderate_entry_weak, moderate_entry_strong
Minimum theorem margin & 0.00241 & $8.01\times10^{-4}$\\ % source: base_minimum_theorem_margin, moderate_minimum_theorem_margin
\bottomrule
\end{tabular}
\end{columns}}
\medskip
\GrayLine{Units are arbitrary; magnitudes are not calibrated; the signs are the theorem; nonemptiness holds for every $h>\ell$ (A17).}
\Takeaway{A moderate economy with $h=2$ % source: moderate_h
gives a rise of similar size (0.250 % source: moderate_entry_weak
$\to$ 0.527), % source: moderate_entry_strong
so the benchmark scale is a declaration, not the mechanism.}
\BackButton{main:scale}
\end{frame}

% ---------------------------------------------------------------- A20
\begin{frame}[label=A20]{How entry and ownership are computed}
\hypertarget{app:entry}{}
\vspace{-2pt}
\begin{gather*}
\tau=\frac{\Cost{c_H}-g_L}{g_H-g_L}\quad\text{(\Cost{$c_H$ = expensive cost})},\qquad x^*=\frac b2\log\frac{\tau}{1-\tau},\qquad \alpha_\theta=\Pr(X\ge x^*\mid\theta)\\[4pt]
e_H=\rho+(1-\rho)\alpha_H,\qquad e_L=\rho+(1-\rho)\alpha_L,\qquad \mathsf E=\frac{e_H+e_L}{2},\qquad \mathsf O_H=\frac{e_H}{2}
\end{gather*}
{\small
\begin{itemize}\setlength{\itemsep}{2pt}
\item $\tau$ = belief threshold for costly preparation; $\alpha_\theta$ = probability that order flow crosses the entry threshold; entry $\mathsf E$ = Pr(challenger prepares); high-value ownership $\mathsf O_H$ = Pr(high-value challenger acquires the target)
\item The high-cost window puts $\tau\in(1/2,M)$ and $x^*\in(0,1)$; $\alpha_H>\alpha_L$
\item Above the ceiling strength $r_C\approx3.59$, % source: base_r_high_cost_ceiling
$\tau>M$; the left-limit entry at $r_C$ is 0.506 % source: base_laplace_entry_ceiling_left_limit
\item Expected target proceeds 0.393 / 0.872 / 0.664 at $r_0$ / $r_1$ / $r_2$ % source: tables/table2_equilibrium_controls.tex Panel A
\end{itemize}}
\BackButton{main:entry}
\end{frame}

% ---------------------------------------------------------------- A21
\begin{frame}[label=A21]{Which equilibrium? What the proof certifies}
\hypertarget{app:cert}{}
\vspace{5pt}
{\small The reversal (frames 9--10) compares economies in which trading and on-path entry are unique, so no selection is needed; multiplicity arises only between them (frame 12).\par}
\vspace{4pt}
\begin{center}\footnotesize
\begin{tabular}{@{}cccc@{}}
\toprule
$r$ & $q_L\in$ & Entry $\mathsf E\in$ & High-type cover margin $\ge$\\
\midrule
1.55 % source: cert_a_r
 & $[-0.46031620,\,-0.46031618]$ % source: cert_a_v_interval (negated)
 & $[0.5450528898,\,0.5450528922]$ % source: cert_a_entry_interval
 & 0.0000761777\\ % source: cert_a_high_derivative_lower
1.60 % source: cert_b_r
 & $[-0.70747539,\,-0.70747537]$ % source: cert_b_v_interval (negated)
 & $[0.5487563062,\,0.5487563085]$ % source: cert_b_entry_interval
 & 0.0027531948\\ % source: cert_b_high_derivative_lower
1.65 % source: cert_c_r
 & $[-0.90333201,\,-0.90333198]$ % source: cert_c_v_interval (negated)
 & $[0.5513607988,\,0.5513608020]$ % source: cert_c_entry_interval
 & 0.0054921767\\ % source: cert_c_high_derivative_lower
\bottomrule
\end{tabular}
\end{center}
\vspace{-4pt}
{\small
\begin{itemize}\setlength{\itemsep}{2pt}
\item \textbf{Low type:} global strict concavity; a sign change of its marginal profit at its own order brackets the root (0.000000000114 % source: cert_a_psi_left_lower
and $-0.000000000096$ at $r=1.55$) % source: cert_a_psi_right_upper
\item \textbf{High type:} cover margin strictly positive over the whole bracket (last column)
\item \textbf{Not proved:} uniqueness of the informative profile, absence of mixed equilibria, a branch between nodes
\end{itemize}}
\Status{computer-assisted (Proposition 3): interval arithmetic on exact decimal inputs}\par
\BackButton{main:cert}
\end{frame}

% ---------------------------------------------------------------- A22
\begin{frame}[label=A22]{Information controls in detail}
\hypertarget{app:controls}{}
{\small
\begin{itemize}\setlength{\itemsep}{0pt}
\item Prop.~A.3 \Status{(analytical)}: at any fixed information experiment and cost law, entry weakly falls with strength; within fixed orders, not across order profiles that change with $r$
\item The frozen weak profile is a control, not an equilibrium: full orders lose money there because $\Delta_T(r_0)<k$; at $r_1$ it coincides with the equilibrium
\end{itemize}}
\vspace{-2pt}
% source: tables/table2_equilibrium_controls.tex Panels A-B; 3 significant digits
\vspace{2pt}
{\scriptsize
\renewcommand{\arraystretch}{0.86}
\setlength{\tabcolsep}{5pt}
\noindent\begin{tabular}{@{}>{\raggedright\arraybackslash}p{4.1cm}cc>{\centering\arraybackslash}p{3.3cm}>{\centering\arraybackslash}p{2.1cm}l@{}}
\toprule
 & $r$ & Entry $\mathsf E$ & High-value ownership $\mathsf O_H$ & Target proceeds & Status\\
\midrule
\emph{Equilibria of the feedback game} & 1.2 & 0.250 & 0.125 & 0.393 & \Status{analytical}\\
 & 3 & 0.523 & 0.324 & 0.872 & \Status{analytical}\\
 & 3.6 & 0.250 & 0.125 & 0.664 & \Status{analytical}\\
\midrule
\multicolumn{6}{@{}l}{\emph{Frozen informative orders} (control, not an equilibrium at $r_0$)}\\
 & 1.2 & 0.562 & 0.351 & 0.520 & \Status{numerical diagnostic}\\
 & 3 & 0.523 & 0.324 & 0.872 & \Status{numerical diagnostic}\\
\midrule
\multicolumn{6}{@{}l}{\emph{Price hidden from challenger} (equilibrium of the no-price-access game)}\\
 & 1.2 & 0.250 & 0.125 & 0.393 & \Status{analytical}\\
 & 3 & 0.250 & 0.125 & 0.615 & \Status{analytical}\\
\bottomrule
\end{tabular}\par}
\GrayLine{Price hidden: the challenger enters only at low cost (high-cost window); at $r_1$ the investor still trades full orders.}
\BackButton{main:controls}
\end{frame}

% ---------------------------------------------------------------- A23
\begin{frame}[label=A23]{Price level or information?}
\hypertarget{app:welfare}{}
\begin{itemize}\setlength{\itemsep}{6pt}
\item \textbf{Invariance diagnostic:} add a dividend $D_0=0.257809$ % source: base_matched_dividend
(the revenue gain from price access) to the traded claim in the price-hidden economy. Mean prices then match, but entry does not, so the information matters, not the price level. \Status{(diagnostic)}
\item The revenue gain itself, 0.258, is analytical % source: base_revenue_gain
\item Prop.~A.9 at $r_1$: each extra entry is chosen only when expected gross profit covers its cost; the allocation gain includes that profit and sales that would otherwise fail the reserve; transfers excluded
\item The dividend is not a sale mechanism and is excluded from surplus; no ranking of strengths or mechanisms
\end{itemize}
\BackButton{main:welfare}
\end{frame}

% ---------------------------------------------------------------- A24
\begin{frame}[label=A24]{Trading and entry across strengths: entry}
\hypertarget{app:corr}{}
\vspace*{-6pt}
\centering
% X5 panels are drawn at 6.0 x 2.42 in with 10 pt text; at 0.9 linewidth the text is 9 pt on the slide
\includegraphics[width=0.9\linewidth]{equilibrium_correspondence_a_talk.pdf}
\par\raggedright
{\footnotesize No trade unique for $r<r_P\approx1.22$, % source: base_r_pool_unique_sufficient
an equilibrium up to $r_N\approx1.75$; % source: base_r_no_trade_exact
full orders unique above $r_U\approx2.84$; % source: base_r_full_unique_sufficient
expensive entry impossible above $r_C\approx3.59$ % source: base_r_high_cost_ceiling
\Status{(analytical, Prop.~A.4; shaded: unique)}; full correspondence open.\par}
\NavGap
\PlaceNav{\hyperlink{app:corrb}{\beamergotobutton{Order sizes}}\NavSep\hyperlink{main:corr}{\beamerreturnbutton{Back}}}
\end{frame}

% ---------------------------------------------------------------- A25
\begin{frame}[label=A25]{Trading and entry across strengths: orders}
\hypertarget{app:corrb}{}
\centering
\includegraphics[width=0.9\linewidth]{equilibrium_correspondence_b_talk.pdf}
\par\raggedright
{\footnotesize No trade is $(0,0)$. Certified nodes: $|q_L|\approx0.460,\,0.707,\,0.903$ % source: cert_a_v_interval, cert_b_v_interval, cert_c_v_interval (negated q_L; absolute values)
(computer-assisted); the curve through them and the mixed candidate are numerical diagnostics; no branch between nodes is claimed.\par}
\BackButton{app:corr}
\end{frame}

% ---------------------------------------------------------------- A26
\begin{frame}[label=A26]{Robustness in full}
\hypertarget{app:robust}{}
\vspace{4pt}
% source: tables/table3_extensions.tex; 3 significant digits, pp as printed
\begin{center}\footnotesize
\setlength{\tabcolsep}{4.5pt}
\begin{tabular}{@{}llccccc@{}}
\toprule
 & $\rho$; strengths & $\mathsf E$ weak & $\mathsf E$ strong & Change (pp) & $\mathsf O_H$ weak & $\mathsf O_H$ strong\\
\midrule
Laplace, cost atoms & 0.25; $1.2\to3$ & 0.250 & 0.523 & 27.28 & 0.125 & 0.324\\ % source: base_rho, base_r_weak, base_r_strong; base_entry_change_pp
Laplace, cost mixture & 0.25; $1.2\to3$ & 0.250 & 0.523 & 27.27 & 0.125 & 0.324\\ % source: tables/table3_extensions.tex; cost_mix_laplace_entry_strong
Logistic, cost atoms & 0.25; $1.2\to3$ & 0.250 & 0.302 & 5.15 & 0.125 & 0.162\\ % source: tables/table3_extensions.tex; logistic_entry_strong
Logistic, cost mixture & 0.25; $1.2\to3$ & 0.250 & 0.301 & 5.14 & 0.125 & 0.162\\ % source: tables/table3_extensions.tex; cost_mix_logistic_entry_strong
Moderate values ($h=2$) & 0.25; $1.05\to1.5$ & 0.250 & 0.527 & 27.68 & 0.125 & 0.327\\ % source: moderate_rho, moderate_r_weak, moderate_r_strong; moderate_entry_*; moderate_entry_change_pp
Signals ($a=0.70$, $d=0.75$) & 0.85; $1.1\to2.3$ & 0.850 & 0.879 & 2.94 & 0.425 & 0.449\\ % source: signal_rho, signal_r_weak, signal_r_strong; signal_entry_*; signal_entry_change_pp
\bottomrule
\end{tabular}
\end{center}
\vspace{-4pt}
{\small
\begin{itemize}\setlength{\itemsep}{2pt}
\item Minimum theorem margins: base $2.41\times10^{-3}$, % source: base_minimum_theorem_margin
moderate $8.01\times10^{-4}$, % source: moderate_minimum_theorem_margin
signals $1.45\times10^{-3}$ % source: signal_minimum_theorem_margin
\item Distinct parameter vectors, not one joint calibration; all rows analytical (Props.~2, A.5--A.7); entry $\mathsf E$ = Pr(challenger prepares), $\mathsf O_H$ = high-value ownership; pp from unrounded probabilities
\end{itemize}}
\Takeaway{The rise from $r_0$ to $r_1$ holds in every row; the $r_2$ fall is shown for the benchmark only.}
\BackButton{main:robust}
\end{frame}

% ---------------------------------------------------------------- A27
\begin{frame}[label=A27]{A higher reserve can raise proceeds; optimal terms are open}
\hypertarget{app:reserve}{}
{\small Value classes: $p=0.5$ % source: base_p
vs $p=1.1$ % source: value_reserve_high
($\varepsilon_V=0.05$)\par} % source: value_band_halfwidth
\vspace{2pt}
% source: tables/table4_reserve_comparisons.tex Panel B; 3 significant digits
\begin{center}\footnotesize
\begin{tabular}{@{}lccccl@{}}
\toprule
Economy, reserve & Preparation & Sale & Two admissible bidders & Expected proceeds & Trading\\
\midrule
Weak $r=1.2$, $p=0.5$ & 0.250 & 0.688 & 0.146 & 0.393 & no trade\\ % source: tables/table4_reserve_comparisons.tex Panel B; value_revenue_weak_low_p
Weak $r=1.2$, $p=1.1$ & 0.540 & 0.392 & 0.0280 & 0.432 & full orders\\ % source: value_entry_weak_high_p; tables/table4_reserve_comparisons.tex Panel B; value_revenue_weak_high_p
Strong $r=3$, $p=0.5$ & 0.523 & 0.920 & 0.436 & 0.872 & full orders\\ % source: tables/table4_reserve_comparisons.tex Panel B; value_revenue_strong_low_p
Strong $r=3$, $p=1.1$ & 0.512 & 0.749 & 0.200 & 1.01 & full orders\\ % source: value_entry_strong_high_p; tables/table4_reserve_comparisons.tex Panel B; value_revenue_strong_high_p
\bottomrule
\end{tabular}
\end{center}
\vspace{2pt}
\GrayLine{Binary values: $p=0.5$ vs $p=1.01$. \Status{Analytical at the listed nodes.}} % source: base_p; tables/table4_reserve_comparisons.tex Panel A
\vspace{4pt}
\begin{itemize}\setlength{\itemsep}{3pt}
\item Preparation, sale and two admissible bidders come apart
\end{itemize}
\TakeawayWithNav{A feasible improvement, not an optimal reserve: optimal terms are open.}
\PlaceNav{\hyperlink{app:pool}{\beamergotobutton{Same orders, different prices}}\NavSep\hyperlink{main:reserve}{\beamerreturnbutton{Back}}}
\end{frame}

% ---------------------------------------------------------------- A28
\begin{frame}[label=A28]{Same orders, different prices}
\hypertarget{app:pool}{}
{\small Prop.~A.10 (analytical existence): at reserve $p=7$ % source: pool_reserve
and $r_0$, the same full orders $(1,-1)$ support price pools below a cutoff $\kappa\in[-\log2,\,0]$.\par} % source: pool_posterior_cutoff_low / pool_posterior_cutoff_high branch labels (c=-log2*(1-0/16), c=-log2*(1-16/16))
\vspace{4pt}
\begin{center}\small
\begin{tabular}{@{}lcc@{}}
\toprule
 & $\kappa=-\log2$ % source: pool_posterior_cutoff_low branch label (c=-log2*(1-0/16))
 & $\kappa=0$\\ % source: pool_posterior_cutoff_high branch label (c=-log2*(1-16/16))
\midrule
Pooled belief & 0.273 % source: pool_posterior_cutoff_low
 & 0.303\\ % source: pool_posterior_cutoff_high
Entry & 0.152 % source: pool_entry_cutoff_low
 & 0.125\\ % source: pool_entry_cutoff_high
Expected target proceeds & 0.687 % source: pool_revenue_cutoff_low
 & 0.610\\ % source: pool_revenue_cutoff_high
\bottomrule
\end{tabular}
\end{center}
\begin{itemize}\setlength{\itemsep}{3pt}
\item Same orders, different prices, beliefs, and entry
\item Does not arise on the benchmark support
\end{itemize}
\Takeaway{A continuation must include the price rule, so the seller's problem cannot be solved over orders alone.}
\BackButton{app:reserve}
\end{frame}

% ---------------------------------------------------------------- A29
\begin{frame}[label=A29]{References}
\hypertarget{app:refs}{}
{\footnotesize
\begin{columns}[T]
\column{0.49\textwidth}
\raggedright
Betton, Eckbo, Thompson and Thorburn (2014). \emph{J. Finance} 69.\par\smallskip
Boone and Mulherin (2007). \emph{J. Finance} 62.\par\smallskip
Bulow, Huang and Klemperer (1999). \emph{J. Polit. Econ.} 107.\par\smallskip
Carlin, Liu, Officer, Pernoud and Tu (2026). NBER Working Paper 34846.\par\smallskip
Cornelli and Li (2002). \emph{Rev. Financ. Stud.} 15.\par\smallskip
Dow, Goldstein and Guembel (2017). \emph{J. Eur. Econ. Assoc.} 15.\par\smallskip
Edmans, Goldstein and Jiang (2012). \emph{J. Finance} 67.\par\smallskip
Edmans, Goldstein and Jiang (2015). \emph{Amer. Econ. Rev.} 105.\par\smallskip
Fishman (1988). \emph{RAND J. Econ.} 19.
\column{0.49\textwidth}
\raggedright
Gentry and Stroup (2019). \emph{J. Financ. Econ.} 132.\par\smallskip
Goldstein and Guembel (2008). \emph{Rev. Econ. Stud.} 75.\par\smallskip
Grossman and Hart (1980). \emph{Bell J. Econ.} 11.\par\smallskip
Hirshleifer and Png (1989). \emph{Rev. Financ. Stud.} 2.\par\smallskip
Imprivata, Inc. (2016). Definitive merger proxy, Form DEFM14A.\par\smallskip
Levin and Smith (1994). \emph{Amer. Econ. Rev.} 84.\par\smallskip
Lin, Ma, Yang and Zhu (2025). Working paper.\par\smallskip
Liu and Bernhardt (2022). Working paper.\par\smallskip
Luo (2005). \emph{J. Finance} 60.\par\smallskip
Pernoud and Gleyze (2026). Working paper.\par\smallskip
Persico (2000). \emph{Econometrica} 68.\par\smallskip
Roberts and Sweeting (2013). \emph{Amer. Econ. Rev.} 103.
\end{columns}}
\BackButton{main:refs}
\end{frame}

\end{document}
